\section{DH description}
\label{sec:dh}

The six revolute joints are described using the \emph{standard} DH convention.
Frame $\frameid{0}$ coincides with \texttt{base\_link}, and $q_i=0$ denotes the
URDF home configuration. The endpoint $E$ is \texttt{link\_6}.
For each step $i$, the transform from frame $\frameid{i-1}$ to frame $\frameid{i}$ is
\begin{equation}
 {}^{i-1}\!T_i
 =R_z(\theta_i)\,T_z(d_i)\,T_x(a_i)\,R_x(\alpha_i),
 \qquad \theta_i=q_i+\theta_i^0 .
 \label{eq:dh-convention}
\end{equation}
Joint $i$ rotates about $z_{i-1}$. For $i<6$, the new axis $z_i$
lies on joint $i+1$; the terminal axis $z_6$ is chosen parallel to $z_5$.
DH frame numbers therefore differ from physical joint numbers.

\textbf{How to read one step.}
Choose $x_i$ along a common normal to $z_{i-1}$ and $z_i$. Its intersection
with $z_i$ defines $O_i$. If two distinct axes intersect, their common-normal length is
zero and $x_i$ is perpendicular to both axes through their intersection.
For coincident axes, the origin and transverse $x$ axis are free; here the last
origin is the endpoint.
Set $y_i=z_i\times x_i$ to complete a right-handed frame. The four parameters
then have the following geometric meaning:
\begin{center}
\small
\begin{tabularx}{\linewidth}{@{}c l X@{}}
\toprule
Parameter & Measured about / along & Construction from $\frameid{i-1}$ to $\frameid{i}$ \\
\midrule
$\theta_i$ & about $z_{i-1}$ & Turn $x_{i-1}$ towards $x_i$. \\
$d_i$ & along $z_{i-1}$ & Move from $O_{i-1}$ to the foot of the common normal. \\
$a_i$ & along $x_i$ & Follow the common normal to $O_i$. \\
$\alpha_i$ & about $x_i$ & Turn $z_{i-1}$ towards $z_i$. \\
\bottomrule
\end{tabularx}
\end{center}
Translations are signed along the indicated positive axes; rotations follow the
right-hand rule. Figure~\ref{fig:dh-steps} shows the construction at $q=0$.

\begin{table}[h!]
\centering
\caption{Six DH steps. Lengths are in metres; angular constants are in radians.}
\label{tab:dh}
\small
\setlength{\tabcolsep}{5pt}
\renewcommand{\arraystretch}{1.35}
\begin{tabularx}{\linewidth}{@{}c c c c c X@{}}
\toprule
$i$ & $d_i$ & $\theta_i$ & $\alpha_i$ & $a_i$ & Geometric construction at home \\
\midrule
1 & $1$ & $q_1$ & $-\pi/3$ & $0.8$
  & Rise $0.5+0.5$ from the base, then move $0.8$ along $x_1$ to joint 2.
    Its axis gives a $-60^\circ$ twist. \\
2 & $1.25$ & $q_2$ & $+\pi/2$ & $2$
  & Move $1.25$ along joint 2's axis and $2$ along $x_2$ to joint 3;
    twist by $+90^\circ$. \\
3 & $0.75$ & $q_3+\pi/2$ & $+\pi/2$ & $0$
  & Extend joints 3 and 4's axes to their intersection $O_3$.
    Turn $x_2$ by $+90^\circ$ to select $x_3$; no common-normal translation is needed. \\
4 & $2.5$ & $q_4$ & $-\pi/2$ & $0$
  & From $O_3$, follow joint 4's axis by $1.5+1$ to the wrist centre $O_4$;
    twist by $-90^\circ$. \\
5 & $0$ & $q_5$ & $+\pi/2$ & $0$
  & Joints 5 and 6's axes meet at the wrist centre.
    Keep $O_5=O_4$ and twist by $+90^\circ$. \\
6 & $1.5$ & $q_6$ & $0$ & $0$
  & Follow joint 6's axis from the wrist centre to the endpoint.
    Choose $z_6=z_5$ and no fixed twist. \\
\bottomrule
\end{tabularx}
\end{table}

\textbf{Frame origins versus joint centres.}
A DH origin lies on the selected joint axis but need not coincide with a URDF
joint origin. In particular, $O_3$ lies $1.5$ m before the
physical centre of joint 4. This distance enters $d_4$, giving $d_4=1.5+1=2.5$ m,
while $a_3=0$. Similarly, $d_5=0$ even though the centres of joints 5 and 6 are
$1.5$ m apart: the two axes already intersect at $O_4=O_5$.

\textbf{Endpoint alignment.}
Frame $\frameid{6}$ and the URDF endpoint have the same origin but different axes:
$x_E=z_6$, $y_E=x_6$ and $z_E=y_6$. Hence the forward kinematics uses the fixed
alignment on the right:
\begin{equation}
 {}^0T_E=\left(\prod_{i=1}^{6}{}^{i-1}\!T_i\right){}^6T_E,
 \qquad
 {}^6T_E=\begin{bmatrix}
 0&1&0&0\\0&0&1&0\\1&0&0&0\\0&0&0&1
 \end{bmatrix}.
 \label{eq:dh-tool}
\end{equation}

\clearpage
\begin{figure}[p]
\centering
{\large\sffamily\bfseries\color{dhink}The six frame-to-frame constructions}\par
\vspace{4pt}
{\small Home configuration ($q=0$), with a fixed viewing direction throughout.}\par
\vspace{7pt}

\DHPanel{1\quad Base to shoulder: $\frameid{0}\to\frameid{1}$}
 {$d_1=1,\quad\theta_1=q_1,\quad\alpha_1=-60^\circ,\quad a_1=0.8$}{1}
 {The full base height is $1$ m. The $30^\circ$ URDF roll of joint 2 translates into a an $\alpha$ of
 $-60^\circ$ in DH.}
\hfill
\DHPanel{2\quad Shoulder to elbow: $\frameid{1}\to\frameid{2}$}
 {$d_2=1.25,\quad\theta_2=q_2,\quad\alpha_2=90^\circ,\quad a_2=2$}{2}
 {The $x$ axes remain parallel. The offset along $z_1$ and the common normal
 along $x_2$ locate $O_2$ on joint 3.}

\vspace{10pt}
\DHPanel{3\quad Axis intersection: $\frameid{2}\to\frameid{3}$}
 {$d_3=0.75,\quad\theta_3=q_3+90^\circ,\quad\alpha_3=90^\circ,\quad a_3=0$}{3}
 {$O_3$ is the intersection of the extended axes. Choosing $x_3=y_2$
 explains the fixed $+90^\circ$ offset in $\theta_3$.}
\hfill
\DHPanel{4\quad Reach the wrist: $\frameid{3}\to\frameid{4}$}
 {$d_4=2.5,\quad\theta_4=q_4,\quad\alpha_4=-90^\circ,\quad a_4=0$}{4}
 {The distance starts at $O_3$, rather than the joint-4 cylinder.
 The two physical offsets add: $d_4=1.5+1$.}

\vspace{10pt}
\DHPanel{5\quad Coincident origins: $\frameid{4}\to\frameid{5}$}
 {$d_5=0,\quad\theta_5=q_5,\quad\alpha_5=90^\circ,\quad a_5=0$}{5}
 {The origins coincide, although the joint centres do not.
 Here $x_4=x_5$ and the extended joint-6 axis passes through the wrist centre.}
\hfill
\DHPanel{6\quad Wrist to endpoint: $\frameid{5}\to\frameid{6}$}
 {$d_6=1.5,\quad\theta_6=q_6,\quad\alpha_6=0,\quad a_6=0$}{6}
 {At home the two frames have parallel axes. $O_6=O_E$;
 the fixed endpoint alignment changes orientation only.}

\vspace{9pt}
\caption{Standard DH frame assignment. Cylinders denote physical revolute joints
 $\jointid{i}$; dots mark DH origins $O_i$.
 \textcolor{dhred}{$x$}, \textcolor{dhgreen}{$y$} and \textcolor{dhblue}{$z$}
 axes retain their colours. Dashed lines extend joint axes;
gold dimensions locate the origins. Angle sketches show $\theta_3^0$ and the signed twist
 $z_{i-1}\to z_i$ about $+x_i$ (with the first direction translated to the same
 point). Cylinder orientation, size and viewing direction are reused between panels;
 drawings are schematic.}
\label{fig:dh-steps}
\end{figure}
